\section{Literature Search and Evidence Assessment}
\label{sec:review-method}

\subsection{Search Scope}

The search began from the target formula in both its structural and its elemental spelling, using \texttt{Ba5Y12Zn[O(SiO4)]8}, \texttt{Ba5Y12ZnSi8O40}, \texttt{BYZSO}, and barium yttrium zinc silicate, combined with the keywords space group, crystal structure, flux, solid-state, phase diagram, and crystal growth. We then followed the literature in three directions: (i) the nomenclature and structural family of the Zn-free tetragonal Ba--Y silicates; (ii) Ba--Zn--Si--O compounds with firm structural evidence; and (iii) Y--Si--O systems for which high-temperature solution growth, solid-state reaction, mechanical activation, or Czochralski processes have been reported. A local full-text materials library was queried first; metadata and experimental passages were then verified through OpenAlex, Crossref, publisher pages, DOI resolution pages, IUCr, COD, RSC main texts and supplementary information (SI), and institutional repositories. The search cutoff was 2026-09-01. A miss in a database or interface means only that that channel returned nothing; it is not evidence that the literature does not exist.

\subsection{Inclusion, Exclusion, and Deduplication}

A publication was included if it met any of the following conditions: it directly reports the preparation or structure of the target phase; it revises the structure, phase region, or compositional series of the tetragonal Ba--Y silicate family; or it provides traceable structural data or synthesis conditions for the structurally related Ba--Zn--Si--O or Y--Si--O systems. Review articles served only to harmonize terminology and to locate primary papers; specific experimental conditions were always taken from the primary studies. We excluded materials related only by element or function, with no crystallographic or synthetic link; phases predicted computationally but never realized experimentally; and aggregated abstracts or product information whose sources could not be traced. Glasses and glass-ceramics were treated as process-reference systems only where they clarify crystallization boundaries or competing phases; they support no inference about the structure of the target phase.

Records were first merged by DOI; where a DOI was missing or several versions existed, they were cross-deduplicated by title, authors, year, and source, and author order, volume and pages, and publication status were checked. The search returned 63 candidate records, of which 51 independent publications survived deduplication and source verification. Because the target phase has only one direct primary paper, we claim neither comprehensive coverage nor a completed independent reproduction.

\subsection{Relevance Classification and Condition Extraction}

We sorted the literature into five classes by how the subject of each study relates to the target phase: direct reports of the target compound; Zn-free members of the same tetragonal Ba--Y silicate family; structurally related compounds with a definite crystallographic relationship; synthesis studies usable as process references; and unrelated sources that cannot support any structural or recipe judgment. The classification applies to individual claims rather than to whole papers; the structural results and the process information from a single source may support conclusions of different scope.

For comparison, we extracted from each publication the starting materials and nominal ratios, pretreatment, mode of phase formation, temperature--time program, atmosphere and whether the system was open or closed, crucible or support, flux, cooling or growth program, and products and characterization methods. Numerical values stay attached to the original experimental record they came from and are never pooled across studies. An item the original text does not state explicitly is recorded as not reported, meaning that this check did not obtain the information, not that the original paper necessarily omitted it. Overall judgments favor results on which several studies agree for the same experimental item, and single-crystal structure, bulk phase composition, measured stoichiometry, and high-temperature stability are assessed separately.
